% !TEX root = ../main.tex
\section{페이지랭크와 순위 상관의 기본 성질}
\label{ch:appendix-theory}

\ref{ch:literature}--\ref{ch:results}장은 여기에 정리한 기본 성질에 기대고 있어요. 이 내용은 널리 알려진 표준적인 것이에요. 페이지랭크의 성질은 \textcite{page1999pagerank}의 연구와 \textcite{langville2006pagerank}의 연구를 따라요. 순위 상관(두 줄 세우기가 얼마나 비슷한지 나타내는 값)은 \textcite{spearman1904proof}의 연구와 \textcite{kendall1938rank}의 연구를 따라요.

\dissertationSubheading[sec:pr-existence]{페이지랭크의 존재와 유일성}

$P$를 $|V|$개의 점 위에 있는 행 확률 행렬(각 행의 수를 모두 더하면 1인 행렬)이라고 해요. 필요하면 식~\eqref{eq:transition}처럼 매달린 점(나가는 화살표가 없는 점)을 고친 뒤의 행렬이에요. $d\in(0,1)$일 때 식~\eqref{eq:google-matrix}의 구글 행렬은 다음과 같아요.
\[
G_{\mathrm{PR}}=d P^{\top}+\frac{1-d}{|V|}\mathbf{1}\mathbf{1}^{\top},
\]
이 행렬은 모든 칸이 양수이고, 열 확률 행렬(각 열의 수를 모두 더하면 1인 행렬)이에요. 양수 행렬에 대한 페론--프로베니우스 정리(Perron--Frobenius theorem)에 따르면, $G_{\mathrm{PR}}$에는 크기(절댓값)가 1인 고윳값(행렬을 곱해도 방향이 바뀌지 않는 벡터가 늘어나는 배수)이 하나뿐이에요. 그리고 그 고윳값에 맞는 고유 벡터 $\mathbf{r}$은 모든 칸이 0보다 커요. 우리는 이 벡터를 $\sum_i r_i=1$이 되도록 크기를 맞출 수 있어요. 따라서 감쇠 계수(화살표를 따라갈 확률)와 매달린 점 규칙을 정해 두면, 방향이 있는 모든 그래프는 페이지랭크 점수를 단 하나만 가져요\parencite{langville2006pagerank}.

이 논문의 재시작 벡터 가운데 세 개는 고르지 않아요. AWP의 재시작 벡터는 활발함에 비례해요. 그래서 아무것도 보내지 않은 주소에서는 0이에요. 활동 재시작(활발한 주소에서 더 자주 다시 출발하기)을 쓰는 엔도스랭크(EndorseRank)의 재시작 벡터는 아무것도 허락하지 않은 주소에서 0이에요. S-PR의 재시작 벡터는 허락 걷기의 점수를 정규화(합이 1이 되게 맞춤)한 것이에요. 그래서 송금 층의 점 가운데 허락 점수가 없는 점에서는 0이에요. 따라서 세 구글 행렬 가운데 어느 것도, 모든 칸이 양수인 행렬이 아니에요. 그래도 고정점(계산을 한 번 더 해도 바뀌지 않는 값)은 여전히 하나뿐이에요. 두 확률 벡터의 차이를 구하면 다시 출발 항이 서로 지워져요. 그리고 이 연산자는 두 벡터 사이의 $\|\cdot\|_1$ 거리를 많아야 $d$배로 만들어요. 하지만 일반적으로 고정점의 모든 칸이 0보다 크지는 않아요. 다시 출발 몫이 없고, 다시 출발하는 점에서 출발한 어떤 길로도 닿지 않는 점은 점수가 0으로 남기 때문이에요. 답이 하나뿐이라는 것은, 매개변수를 정해 두면 각 방법이 주어진 화살표 집합에 점수 벡터를 꼭 하나만 준다는 뜻이에요. 그래서 두 점수 벡터 $\mathbf{s}_{\mathrm{ER}}$, $\mathbf{s}_{\mathrm{AWP}}$ 사이의 차이는 모두 두 방법의 화살표와 재시작 분포(어디서 다시 출발하는지의 비율)에서 생겨요.

\dissertationSubheading[sec:pr-continuity]{화살표 무게에 대한 연속성}

$d\in(0,1)$을 고정하면, 페이지랭크는 $P$의 칸 값에 따라 연속적으로(끊기지 않고 조금씩) 변해요. 같은 점들 위에 두 번째 행 확률 행렬 $\tilde P$가 있고, 그 페이지랭크 벡터를 $\tilde{\mathbf{r}}$로 써요. 그러면 다음이 성립해요.
\[
\|\mathbf{r}-\tilde{\mathbf{r}}\|_1\le\frac{d}{1-d}\,\max_i\bigl\|P_{i\cdot}-\tilde P_{i\cdot}\bigr\|_1,
\]
이 상한(넘지 않는 값)은 $|V|$에 따라 달라지지 않아요\parencite{langville2006pagerank}.

$P$에 대한 연속성은 화살표 무게에 대한 연속성과 같지 않아요. 행 $P_{i\cdot}=w_{i\cdot}/o_i$는 나가는 무게의 합 $o_i$가 0보다 큰 동안에만 무게에 따라 연속적으로 움직여요. 어떤 점이 첫 번째 나가는 화살표를 얻거나 마지막 나가는 화살표를 잃으면, 그 점의 행은 갑자기 뛰어요. 고르게 퍼진 매달린 점의 행과, 화살표가 가리키는 이웃들에 몰린 행 사이를 오가요. 관련된 무게가 아무리 작아도 그래요. \ref{sec:pagerank-worked-example}절의 계산 예에서 무게가 $10^{-12}$인 화살표 $D\to A$를 더하면, $\mathbf{r}$이 $\|\cdot\|_1$ 노름(차이의 절댓값을 모두 더한 크기)으로 $0.37$만큼 움직여요. 그리고 $r_A$는 $0.1375$에서 $0.3202$로 올라가요.

감쇠 계수를 여러 값으로 바꿔 보는 실험은 연속성으로 설명할 수 있어요. 이 실험은 $P$를 그대로 두고 $d$만 매끄럽게 바꿔요. 그리고 $d<1$이면 페이지랭크 벡터는 $d$에 따라 연속적으로 변해요. 금액 기울기 점검은 일부만 설명할 수 있어요. 모든 $z>0$에 대해 $V(z)>0$이에요. 그래서 기울기 $b$를 바꿔도, 정확한 계산에서는 화살표가 하나도 더해지거나 빠지지 않아요. 하지만 컴퓨터의 배정밀도 계산에서는 아주 작은 금액만 실은 화살표의 무게가 반올림되어 0이 될 수 있어요. 그리고 $b$를 10배 바꾸는 것은 너무 큰 변화라서, 위의 상한은 이런 변화에 대해 별로 말해 주지 못해요. 그래서 그 점검이 안정적인지는 자료에 달려 있어요(표~\ref{tab:robustness-slope}). 실제 자료에서는 $d\in\{0.75,0.85,0.95\}$에 걸쳐 엔도스랭크의 지표 가족(비슷한 지표 묶음) 평균이 많아야 $0.0031$만큼, AWP의 지표 가족 평균이 많아야 $0.021$만큼 움직였어요(표~\ref{tab:robustness-damping}).

\dissertationSubheading[sec:pr-sparsity]{희소성과 반복 한 바퀴의 비용}

0이 아닌 화살표가 $m=|E|$개, 점이 $n=|V|$개 있으면, 듬성듬성한 곱 $P^{\top}\mathbf{x}$에는 $O(m)$번의 산술 연산이 들어요. 반복 한 바퀴마다 길이가 $n$인 벡터들도 새로 고쳐요. 화살표 목록으로 듬성듬성한 연산자를 만드는 일은 $O(m)$이에요. 그래서 $I$번 반복하는 거듭제곱 반복(같은 계산을 되풀이하기)은 $O(m+(m+n)I)$예요. 두 점수는 같은 풀이 프로그램을 써요. 실제로 잰 호출에서 AWP의 반복 한 바퀴에는 엔도스랭크의 반복 한 바퀴보다 약 $1.95$배의 시간이 걸려요. 이것은 화살표 수의 비율인 $2.21$보다 작아요. 그리고 AWP는 반복을 더 많이 해야 해요(\ref{sec:benchmark-results}절). 그래서 $m$도 $I$도 혼자서는 시간 비율을 정하지 못해요.

\dissertationSubheading[sec:spearman-properties]{스피어먼의 \texorpdfstring{$\rho$}{rho}가 가진 성질}

스피어먼의 $\rho$는 $[-1,1]$ 안에 있어요. $\rho=1$이면 완벽한 단조 일치(한쪽이 커지면 다른 쪽도 늘 커짐)를 뜻해요. $\rho=-1$이면 완벽하게 거꾸로인 일치를 뜻해요. 평가 처리 과정은 동점을 제대로 다루는 표준 구현으로 $\rho$와 $\tau_b$를 계산해요. $n$이 크면 $\rho$와 $\tau$는 보통 부호가 같지만, 크기는 달라요. 엔도스랭크나 AWP를 대리 지표 하나와 짝지은 14개 쌍에서(표~\ref{tab:alignment-transfer}부터~\ref{tab:alignment-sybil-stability}까지), 모든 $|\tau|$는 $0.05$를 넘어요. 비율 $\rho/\tau$는 약 $1.2$와 $1.5$ 사이에 있어요. 다만 들어오는 상대 비율에 대한 AWP는 예외로, $\tau=0.055$이고 $\rho=0.163$이에요. 그리고 둘 가운데 늘 $|\rho|$가 더 커요. 계수의 부호가 서로 반대인 쌍은 하나도 없어요.

\dissertationSubheading[sec:kendall-properties]{켄달의 \texorpdfstring{$\tau$}{tau}가 가진 성질}

이 논문의 점수들에서는 동점이 드물어요. 매칭 코호트(모든 점수가 똑같이 점수를 매기는 계약 묶음)에서 어떤 점수도 계약 쌍의 $3.4\%$보다 많이 동점을 만들지 않아요. 엔도스랭크와 AWP의 동점 비율은 \ref{sec:rank-divergence}절에 있어요. 식~\eqref{eq:kendall}에서 $\tau_b$의 분모는 동점 때문에만 $n_0$과 달라요. 그래서 엔도스랭크와 AWP를 비교할 때 이 분모는 $n_0$과 2퍼센트 미만으로 달라요. 그리고 \ref{sec:rank-divergence}절의 쌍 개수로 계산한 $\tau_a$는 보고한 $\tau_b$와 2퍼센트 안쪽으로만 달라요.

대리 지표에 있는 동점도 $\tau_b$에 천장(얻을 수 있는 가장 큰 값)을 만들어요. 들어오는 화살표 수(입차수)나 받은 허락 수(허락 입차수) 같은 정수 대리 지표는 많은 쌍을 동점으로 만들어요. 그래서 그런 대리 지표에 대해 어떤 점수가 다다를 수 있는 가장 큰 $\tau_b$는 동점 쌍의 비율이 클수록 낮아져요. 페이지랭크 점수들은 동점 쌍이 적어요. 그래서 한 대리 지표에서 두 페이지랭크 점수를 비교하면, 이 천장은 둘 사이에서 거의 다르지 않아요. 따라서 그런 대비는 이 천장 때문에 비뚤어지지 않아요. 하지만 페이지랭크와 원시 차수(화살표를 그냥 센 수) 사이의 대비나, 두 대리 지표 사이의 대비는 비뚤어질 수 있어요. 둘의 동점이 서로 다른 만큼 그래요(\ref{sec:statistical-tests}절).
